\documentclass[lang=cn,newtx,10pt,scheme=chinese,thmcnt=section]{elegantbook}

\title{Advanced Alegebra Notes}
\author{Dengyu Li}
\date{\today}
\version{1.0}

\setcounter{tocdepth}{3}

\usepackage{array}
\newcommand{\ccr}[1]{\makecell{{\color{#1}\rule{1cm}{1cm}}}}

\usepackage{cprotect}

%% personalized settings
\usepackage{tocloft}

\renewcommand{\cftsecpresnum}{\S\ }   % 目录中 section 编号前加 §
\renewcommand{\cftsecaftersnum}{}     % 编号后不额外加东西
\addtolength{\cftsecnumwidth}{1.5em}  % 给前面的符号留出宽度

\numberwithin{equation}{section}
\usepackage{empheq}

\newcounter{propertycnt}[section]
\renewcommand{\thepropertycnt}{\thesection.\arabic{propertycnt}}

\newenvironment{numberedproperty}{
  \refstepcounter{propertycnt}
  \par\noindent\textbf{\color{third}性质~\thepropertycnt}\quad\citshape
}{
  \par
}

\begin{document}

\maketitle
\frontmatter

\tableofcontents

\mainmatter

\chapter{行列式}

\section{二阶行列式}
设有二元一次方程组：
\begin{equation}
\begin{cases}
a_{11}x_1 + a_{12}x_2 = b_1 \\
a_{21}x_1 + a_{22}x_2 = b_2
\end{cases}
\end{equation}
用$a_{22}$乘以第一式的两边，用$a_{12}$乘以第二式的两边，然后相减，得：
\[
(a_{11}a_{22}-a_{12}a_{21})x_1 = b_1a_{22}-b_2a_{12}
\]
于是
\[
x_1 = \frac{b_1a_{22}-b_2a_{12}}{a_{11}a_{22}-a_{12}a_{21}}
\]
用类似的办法消去$x_1$，得：
\[ 
x_2 = \frac{b_2a_{11}-b_1a_{21}}{a_{11}a_{22}-a_{12}a_{21}}
\]
注意到分母$a_{11}a_{22}-a_{12}a_{21}$是一个重要的量，它在行列式中有着重要的作用。于是我们引入二阶行列式如下：
\[  
\begin{vmatrix}
a_{11} & a_{12} \\
a_{21} & a_{22}
\end{vmatrix} = a_{11}a_{22}-a_{12}a_{21}
\]
则上述解可以用行列式表示：
\begin{equation}
x_1 = \frac{\begin{vmatrix}b_1 & a_{12} \\
b_2 & a_{22}\end{vmatrix}}{\begin{vmatrix}a_{11} & a_{12} \\ a_{21} & a_{22}\end{vmatrix}}, \quad
x_2 = \frac{\begin{vmatrix}a_{11} & b_1 \\ a_{21} & b_2\end{vmatrix}}{\begin{vmatrix}a_{11} & a_{12} \\ a_{21} & a_{22}\end{vmatrix}}
\end{equation}
\begin{definition}
    二阶行列式的定义为：
    \[|A| = 
    \begin{vmatrix}
    a_{11} & a_{12} \\
    a_{21} & a_{22}
    \end{vmatrix} = a_{11}a_{22}-a_{12}a_{21}
    \]

\end{definition}
\newpage

\section{三阶行列式}
设有三元一次方程组：
\begin{empheq}[left=\empheqlbrace]{align}
a_{11}x_1 + a_{12}x_2 + a_{13}x_3 = b_1 \\
a_{21}x_1 + a_{22}x_2 + a_{23}x_3 = b_2 \\
a_{31}x_1 + a_{32}x_2 + a_{33}x_3 = b_3
\end{empheq}
我们采用待定系数消元法，即对$(1.2.3)$乘以$u$，$(2.2.3)$乘以$v$，$(3.2.3)$乘以$w$，然后三式相加，消去$x_2$和$x_3$，得到：
\begin{empheq}[left=\empheqlbrace]{align}
&(a_{11}u + a_{21}v + a_{31}w)x_1 = b_1u + b_2v + b_3w \\
&a_{12}u + a_{22}v + a_{32}w = 0 \\
&a_{13}u + a_{23}v + a_{33}w = 0
\end{empheq}
对于$(1.2.5)$和$(1.2.6)$,我们把$a_{32}w$和$a_{33}w$移到右边，再除以w得到关于$\dfrac{u}{w}$和$\dfrac{v}{w}$的二元一次方程组：
\begin{equation}
\begin{cases}
a_{12}\dfrac{u}{w} + a_{22}\dfrac{v}{w} = -a_{32} \\
a_{13}\dfrac{u}{w} + a_{23}\dfrac{v}{w} = -a_{33}
\end{cases}
\end{equation}
由1.1节的结论可知，(1.2.7)的解为：
\[
\dfrac{u}{w} = \frac{\begin{vmatrix}-a_{32} & a_{22} \\ -a_{33} & a_{23}\end{vmatrix}}{\begin{vmatrix}a_{12} & a_{22} \\ a_{13} & a_{23}\end{vmatrix}}, \quad
\dfrac{v}{w} = \frac{\begin{vmatrix}a_{12} & -a_{32} \\ a_{13} & -a_{33}\end{vmatrix}}{\begin{vmatrix}a_{12} & a_{22} \\ a_{13} & a_{23}\end{vmatrix}}
\]
令$u=$$\begin{vmatrix}a_{22} & a_{23} \\ a_{32} & a_{33}\end{vmatrix}$, $v=-$ $\begin{vmatrix}a_{12} & a_{13} \\ a_{32} & a_{33}\end{vmatrix}$, $w=$ $\begin{vmatrix}a_{12} & a_{13} \\ a_{22} & a_{23}\end{vmatrix}$,再由1.1节二阶行列式的定义可知：
\[
x_1=\frac{b_1u + b_2v + b_3w}{a_{11}u + a_{21}v + a_{31}w} = \frac{\begin{vmatrix}b_1 & a_{12} & a_{13} \\ b_2 & a_{22} & a_{23} \\ b_3 & a_{32} & a_{33}\end{vmatrix}}{\begin{vmatrix}a_{11} & a_{12} & a_{13} \\ a_{21} & a_{22} & a_{23} \\ a_{31} & a_{32} & a_{33}\end{vmatrix}}
\]
\[
\begin{vmatrix}
    a_{11} & a_{12} & a_{13} \\ a_{21} & a_{22} & a_{23} \\ a_{31} & a_{32} & a_{33} 
\end{vmatrix} = a_{11}\begin{vmatrix}a_{22} & a_{23} \\ a_{32} & a_{33}\end{vmatrix} - a_{12}\begin{vmatrix}a_{21} & a_{23} \\ a_{31} & a_{33}\end{vmatrix} + a_{13}\begin{vmatrix}a_{21} & a_{22} \\ a_{31} & a_{32}\end{vmatrix}
\]
\begin{definition}
    三阶行列式的递归定义为：
    \[|A| = 
    \begin{vmatrix}
    a_{11} & a_{12} & a_{13} \\
    a_{21} & a_{22} & a_{23} \\
    a_{31} & a_{32} & a_{33}
    \end{vmatrix}
    = a_{11}\begin{vmatrix}a_{22} & a_{23} \\ a_{32} & a_{33}\end{vmatrix} - a_{12}\begin{vmatrix}a_{21} & a_{23} \\ a_{31} & a_{33}\end{vmatrix} + a_{13}\begin{vmatrix}a_{21} & a_{22} \\ a_{31} & a_{32}\end{vmatrix}
    \]
\end{definition}
\newpage

\section{$n$阶行列式}
    由两条竖线围成的n行n列的元素组成的式子(数值)称为n阶行列式，记为：
    \begin{equation}
    |A| = 
    \begin{vmatrix}
    a_{11} & a_{12} & \cdots & a_{1n} \\
    a_{21} & a_{22} & \cdots & a_{2n} \\
    \vdots & \vdots & \ddots & \vdots \\
    a_{n1} & a_{n2} & \cdots & a_{nn}
    \end{vmatrix}
    \end{equation}
    称第i行为：$a_{i1}, a_{i2}, \cdots, a_{in}$，第j列为：$a_{1j}, a_{2j}, \cdots, a_{nj}$，称$a_{ij}$为行列式的元素。\\
    元素$a_{11}, a_{22}, \cdots, a_{nn}$称为行列式的主对角线元素，$a_{1n}, a_{2,n-1}, \cdots, a_{n1}$称为行列式的副对角线元素。
    \begin{definition}
        定义元素$a_{ij}$的余子式为$M_{ij}$为由行列式$|A|$中去掉第$i$行和第$j$列后得到的$(n-1)$阶行列式：
        \[M_{ij} =\begin{vmatrix}
            a_{11} & \cdots & a_{1,j-1} & a_{1,j+1} & \cdots & a_{1n} \\
            \vdots & \ddots & \vdots & \vdots & \ddots & \vdots \\
            a_{i-1,1} & \cdots & a_{i-1,j-1} & a_{i-1,j+1} & \cdots & a_{i-1,n} \\
            a_{i+1,1} & \cdots & a_{i+1,j-1} & a_{i+1,j+1} & \cdots & a_{i+1,n} \\
            \vdots & \ddots & \vdots & \vdots & \ddots & \vdots \\
            a_{n1} & \cdots & a_{n,j-1} & a_{n,j+1} & \cdots & a_{nn}
        \end{vmatrix}\]
    \end{definition}
    \begin{definition}
        当$n=1$时，$(1.3.1)$的值为$|A| = a_{11}$，现假设对$n-1$阶行列式已经定义了它们的值，则对$\forall i,j,M_{ij}$的值已经定义，定义n阶行列式$|A|$的值为：
        \begin{equation}
        |A| = a_{11}M_{11} - a_{12}M_{12} + \cdots + (-1)^{1+n}a_{1n}M_{1n}
        \end{equation}
    \end{definition}
    对于任一自然数$n$，$(1.3.2)$式给出了一个计算$n$阶行列式的方法：将$n$阶行列式展开为$n$个$(n-1)$阶行列式的代数和，再化$n-1$阶行列式为$n-2$阶行列式，以此类推，直到计算出$|A|$行列式的值。
    为了使$(1.3.2)$式更为简洁，我们引进代数余子式的概念
    \begin{definition}
    定义元素$a_{ij}$的代数余子式为：
    \[
        A_{ij} = (-1)^{i+j}M_{ij}
    \]
    其中$M_{ij}$为$a_{ij}$的余子式\\
    用代数余子式，$(1.3.2)$式可以写为：
    \begin{equation}
        |A| = a_{11}A_{11} + a_{12}A_{12} + \cdots + a_{1n}A_{1n}
    \end{equation}
    \end{definition}
    \begin{numberedproperty}
        对于任意的$n$阶行列式$|A|$，且：
        \[
            |A_1| = \begin{vmatrix}
                a_{11} & a_{12} & \cdots & a_{1n} \\
                0 & a_{22} & \cdots & a_{2n} \\
                \vdots & \vdots & \ddots & \vdots \\
                0 & 0 & \cdots & a_{nn}
            \end{vmatrix},
            \text{或}\quad
            |A_2| = \begin{vmatrix}
                a_{11} & 0 & \cdots & 0 \\
                a_{12} & a_{22} & \cdots & a_{n2} \\
                \vdots & \vdots & \ddots & \vdots \\
                a_{1n} & a_{2n} & \cdots & a_{nn}
            \end{vmatrix}
        \]
        则$|A_1|,|A_2| = a_{11}a_{22}\cdots a_{nn}$，即行列式的值等于主对角线元素的乘积。
    \end{numberedproperty}

    \begin{proof}
        我们称上式中左边的行列式为上三角行列式(这时$a_{ij} = 0 , \forall i > j$),称右边的行列式为下三角行列式(这时$a_{ij} = 0 , \forall j > i$)，现在用数学归纳法分别证明上述结论。\\
        First(上三角行列式):\\
        当n=1时，结论显然成立。假设上述结论对$n-1$阶行列式成立，既证$n$阶行列式的情形。\\
        由定义1.3.2可知：
        \[|A_1|=a_{11}M_{11}\]
        其中$M_{11}$为$n-1$阶上三角行列式，由归纳假设可知：
        \[M_{11} = a_{22}a_{33}\cdots a_{nn}\]
        因此：
        \[|A_1| = a_{11}a_{22}a_{33}\cdots a_{nn}\]证毕。\\
        Second(下三角行列式):\\
        由定义$1.3.2$可知：
        \[|A_2| = a_{11}M_{11} - a_{12}M_{12} + \cdots + (-1)^{1+n}a_{1n}M_{1n}\]
        考虑余子式$M_{k1}(1 \le k \le n)$,设$M_{k1}=(b_{ij})$其为$n-1$阶行列式，则
        \[
            b_{ij} = 
            \begin{cases}
                a_{i,j+1}, & 1 \le i \le k-1\\
                a_{i+1,j+1}, & k \le i \le n
            \end{cases}
        \]
        由下三角行列式的定义：对于$\forall i > j, a_{ij} = 0$。而对于$b_{ij}$,无论$k$取何值，其都满足$\forall i > j, b_{ij} = a_{i(j+1)} $ 或 $ a_{(i+1)(j+1)}= 0$，因此$M_{k1}$为$n-1$阶下三角行列式，于是有下列断言。\\
        
        \begin{proposition}
            当$k \ge 2$时，$M_{k1}$总有一个主对角线元素为0，其递归展开后的值也必为0。
        \end{proposition}
        \begin{proof}
            当$k \ge 2$时, $b_{\,(k-1)(k-1)} = a_{(k-1)k} $，而由下三角行列式的定义，$a_{(k-1)k} = 0$一定成立，因此$M_{k1}$总有一个主对角线元素为0，由定义$1.3.2$可知，故其递归展开后值必为0
        \end{proof}
        由归纳假设可知，$M_{k1} = 0( k \ge 2)$，$M_{11}= a_{22}a_{33}\cdots a_{nn}$，因此：$|A_2| = a_{11}M_{11} = a_{11}a_{22}a_{33}\cdots a_{nn}$，证毕。\\


        


    \end{proof}

    \begin{numberedproperty}\label{pro:det-one}
        
    \end{numberedproperty}

\chapter{矩阵}

\end{document}